\section{Injection across context tools and persistent memory}\label{sec:C28:injection}
Prompt injection exploits confusion between data and instructions across a trust boundary. An agent reads external content needed for a legitimate task, but that content attempts to alter goals, tool use, or disclosure. The vulnerability is architectural because the same language interface can contain both evidence and directives. A warning in the prompt may help behavior but does not create a hard separation of authority.

Tool outputs can carry several kinds of risk. They may contain untrusted prose, malformed structured fields, misleading status, or references to external artifacts. A tool can be honest while returning attacker-controlled content, so classifying the whole tool as trusted is too coarse. Preserve provenance and distinguish the tool's transport or schema guarantees from the authority of the data it retrieves.

Persistent memory extends the attack horizon. A malicious instruction can be extracted as a preference, consolidated into a summary, or promoted into a reusable skill. Later retrieval can make it appear to be established internal knowledge. The attacker no longer needs to be present in the current task. A memory writer therefore needs provenance, scope and promotion controls, and the reader must revalidate authority rather than trust the storage location.

Skill supply chains add executable and procedural risk. A skill can contain instructions, scripts, dependencies and tool references. Reviewing only its short description misses transitive behavior. Loading a skill should not automatically grant the permissions needed by its embedded resources. Version pinning and application-enforced capabilities limit the effect boundary, while semantic review examines whether the procedure remains applicable to the current task.

MCP and other tool protocols provide structured interfaces, but protocol conformance does not establish tool honesty or policy compliance~\cite{mcp2026spec}. Discovery metadata, annotations and returned content have distinct trust roles. A client must enforce authorization and validate arguments independently of descriptive claims supplied by a server. Schema-valid invocation is a necessary interface property, not a universal security guarantee.

SkillJack makes cross-layer promotion concrete. Its attacker injects experiences without direct write access to the skill library; the unchanged extraction and routing pipeline can turn those records into separately stored procedural artifacts~\cite{ying2026skilljack}. Deleting the original record can then leave the derived skill available. The report's scope is two systems with one model and a policy-violation proxy, not demonstrated live-service compromise. A diagnostic must therefore test source deletion and derivative invalidation separately, retain clean queries to detect unintended activation, and judge effects independently of whether a skill is routed. Changing extraction can reduce recognizable malicious wording without removing the harmful procedural implication.

MemEvoBench probes a different mechanism: repeated mixed correct and misleading memory, biased feedback, and noisy workflow returns can shift later behavior~\cite{xie2026memevo}. Its QA and workflow constructions use generated memory pools; their contamination rate and update protocol define an experimental exposure, not its prevalence in deployment. To identify accumulation rather than immediate persuasion, compare the same final visible evidence with and without intervening memory writes, freeze the reader policy, and counterbalance exposure order. A source-quarantine defense can prevent an initial write yet fail after an unsafe abstraction has been promoted elsewhere. These cases require the dependency-aware revocation and longitudinal units of Chapters~\ref{ch:C13} and~\ref{ch:C27}.

